\chapter{Differentials proper}

\begin{orient}
\textbf{What this chapter is for.} A differential is the linear part of an increment, and almost every confusion about $\dd x$ dissolves the moment that sentence is taken literally. If $y=f(x)$ and $x$ is nudged by an amount $\Delta x$, the true change $\Delta y=f(x+\Delta x)-f(x)$ is some complicated thing; the differential $\dd y=f'(x)\dd x$ is the best straight-line stand-in for it, and the two agree to first order and disagree at second order. Everything in this chapter is a consequence: linear approximation is that stand-in used as an estimate, error propagation is that stand-in used as a tolerance, the total differential is that stand-in in several variables, and $\dd u=g'(x)\dd x$ under a substitution is that stand-in used as a change of coordinates on the increment. By the end you should be able to produce a numerical estimate with a signed error bound, convert a measurement tolerance into a tolerance on any derived quantity, differentiate an implicit relation without ever solving it, and say precisely which manipulations of $\dd x$ are legitimate and why. That last point is the one most textbooks skip, and it is the reason students distrust the notation for years.
\end{orient}

\section{The increment and its linear part}

Fix a differentiable $f$ and a point $x$. The definition of the derivative says
\[f(x+\Delta x)-f(x)=f'(x)\,\Delta x+\varepsilon(\Delta x)\,\Delta x,\qquad \varepsilon(\Delta x)\to0 \text{ as } \Delta x\to0.\]
Read the right side as two objects rather than one. The first term is linear in $\Delta x$ with a coefficient that does not depend on $\Delta x$; the second is a remainder that vanishes faster than $\Delta x$ does. We give the first term a name and a notation: for an independent variable we \emph{define} $\dd x=\Delta x$, and for the dependent variable we \emph{define}
\[\dd y=f'(x)\dd x.\]
Two things follow immediately and are worth saying out loud. First, $\dd y$ is a function of two independent inputs, the base point $x$ and the nudge $\dd x$; it is not a number attached to $x$ alone. Second, $\dd y$ is exactly computable while $\Delta y$ usually is not, which is the entire practical point: we trade an unknown quantity for a linear one and pay a second-order price.

The ratio $\dd y/\dd x$ is then honestly a quotient of two real numbers, and it equals $f'(x)$ for every nonzero $\dd x$. This is why Leibniz notation behaves so well under the chain rule and under substitution, and it is also the precise sense in which $\dd y/\dd x$ ``is'' a fraction. What is \emph{not} legitimate is treating $\dd x$ as an infinitesimal, a number smaller than every positive real. No such real number exists. Everything in this chapter can be done with $\dd x$ an ordinary finite real, and everything that looks like cancelling differentials is really the chain rule or the substitution rule written compactly.

The picture is worth fixing in mind, because every technique below is a reading of it. Draw the graph of $f$ and the tangent line at $x$. Step right by $\dd x$. The graph rises by $\Delta y$; the tangent line rises by $\dd y$. The gap between them is the whole error, it is second order in $\dd x$, and its sign is the sign of $f''$: for a concave function the tangent lies above the curve and $\dd y>\Delta y$, for a convex function the reverse. Every estimate in this chapter is the height of that tangent line, every error bound is the size of that gap, and every question about whether an estimate is high or low is answered by looking at the concavity rather than by computing anything. Halving $\dd x$ halves $\dd y$ and quarters the gap, which is the practical content of ``second order'' and the reason linear approximation is useful at all.

\begin{keynote}[What the notation is, honestly]
There are three defensible readings of $\dd x$, and each licenses a different set of moves. (i) \textbf{Finite increment}: $\dd x$ is any real number, $\dd y=f'(x)\dd x$ is a linear function of it. This is the reading used for approximation and error propagation, and it makes every equation in this chapter a true statement about real numbers. (ii) \textbf{Formal symbol under an integral or in an ODE}: $\dd x$ is a bookkeeping token whose only rule is that $u=g(x)$ forces $\dd u=g'(x)\dd x$, justified by the substitution rule. This is the reading used when you separate variables. (iii) \textbf{Differential form}: $\dd x$ and $\dd y$ are covectors and $\dd F=F_x\dd x+F_y\dd y$ is an exact $1$-form. This is the reading that makes exactness of $M\dd x+N\dd y=0$ a theorem rather than a coincidence. All three agree on every computation below. None of them permits ``multiply both sides by $\dd x$'' as a self-standing act; that phrase is always shorthand for one of (i), (ii), or (iii).
\end{keynote}

One structural property of the notation deserves to be stated because it is what makes the whole scheme coherent, and because it is invisible if you only ever work in one variable. The expression $\dd z=f_x\dd x+f_y\dd y$ is true whether $x$ and $y$ are independent variables or are themselves functions of some other variables $s$ and $t$. In the second case $\dd x$ and $\dd y$ are no longer free but are given by their own differentials, and substituting them reproduces exactly the multivariable chain rule. This is called the invariance of the form of the first differential, and it is the reason you may write $\dd$ of an equation without first deciding which symbols are independent. Every implicit differentiation in this chapter relies on it, and so does every substitution in the next two.

The first technique is nothing but the definition, applied on purpose. It looks too simple to need a name, and it is the technique students most reliably misuse, because they report $\dd y$ where $\Delta y$ was asked for and never notice the difference.

\tech[Differential dy versus increment Delta y]{The differential $\dd y$ versus the increment $\Delta y$}{core}
\cue A problem asks approximately how much $y$ changes when $x$ changes by a small stated amount, or supplies $\Delta x$ and wants $\Delta y$ without recomputing $f$. The word ``approximately'', or a tolerance sign $\pm$, is the tell.
\method Set $\dd x=\Delta x$ and compute $\dd y=f'(x)\dd x$. Then
\[\Delta y=\dd y+o(\dd x),\qquad\text{so}\qquad f(x+\Delta x)\approx f(x)+f'(x)\,\Delta x.\]
The differential is the linear part of the increment, not the increment. Report it as an approximation, and if the question wants a bound, get the sign of the error from the sign of $f''$ near $x$.

\begin{worked}
Estimate $\sqrt{100.4}$. Take $f(x)=\sqrt x$, $x=100$, $\dd x=0.4$. Then $f'(100)=\tfrac{1}{2\sqrt{100}}=\tfrac1{20}$, so
\[\dd y=\frac{0.4}{20}=0.02,\qquad \sqrt{100.4}\approx10.02.\]
The true value is $10.019980\ldots$, so the differential overshot by $2.0\times10^{-5}$. That sign was predictable: $f''(x)=-\tfrac14x^{-3/2}<0$, so $\sqrt{\cdot}$ is concave and its tangent line lies above the curve everywhere. The size was predictable too, from $\tfrac12\abs{f''}(\dd x)^{2}=\tfrac12\cdot\tfrac{1}{4000}\cdot0.16=2.0\times10^{-5}$.

A second instance, where $\Delta y$ and $\dd y$ visibly part company: $y=x^{10}$ at $x=1$ with $\Delta x=0.02$. Here $\dd y=10(0.02)=0.2$, giving the estimate $1.2$, while $(1.02)^{10}=1.218994\ldots$. The differential captures $0.2$ of the true increment $0.218994$; the missing $0.019$ is the quadratic term $\tfrac12\cdot90\cdot(0.02)^{2}=0.018$ plus higher order.
\end{worked}

\begin{keynote}[Rates are differentials divided by $\dd t$]
The related-rate problems of elementary calculus are this technique with both variables made functions of time. If $V=\tfrac43\pi r^{3}$ then $\dd V=4\pi r^{2}\dd r$, and dividing by $\dd t$ gives $\dv{V}{t}=4\pi r^{2}\dv{r}{t}$: the same identity, read as a relation between rates rather than between increments. Nothing new is needed, and setting these problems up in differential form first has a practical advantage, namely that you do not have to decide in advance which variable is independent. Take the differential of the geometric relation, then divide by $\dd t$, and the chain rule appears automatically.
\end{keynote}

\begin{trap}
Reporting $\dd y$ as though it were exact, and forgetting that $\dd y$ depends on two variables. Students silently fix $\dd x=1$, which turns $\dd y$ into $f'(x)$ and destroys the units of the answer: if $x$ is measured in centimetres and $\dd x=0.4$ cm, then $\dd y$ carries the units of $f'\cdot\text{cm}$, and dropping the $0.4$ reports a rate where a quantity was wanted.
\end{trap}

\begin{seealsob}Linear approximation and the tangent-line estimate; error propagation; Taylor polynomials, of which this is the degree-one case.\end{seealsob}

Written as a statement about $f$ rather than about the increment, the same identity is the tangent line, and the shift of viewpoint is worth making because it tells you what to optimise. In the increment form the free choice is $\Delta x$, which the problem usually fixes. In the tangent form the free choice is the base point $a$, which the problem almost never fixes, and choosing $a$ badly is how an otherwise correct estimate becomes useless.

\tech{Linear approximation and the tangent-line estimate}{easy}
\cue ``Estimate $f(a+h)$'' where some nearby point has exactly computable $f$ and $f'$, or the explicit instruction ``use differentials to approximate''. Also any request for a first-order correction to a known value.
\method Use the linearisation
\[L(x)=f(a)+f'(a)(x-a),\qquad f(x)-L(x)=\tfrac12f''(\xi)(x-a)^{2}\ \text{ for some }\xi\text{ between }a\text{ and }x.\]
Choose $a$ to be the \emph{nearest} point at which $f$ and $f'$ are both exact, not the most arithmetically convenient one. The error formula then does two jobs: it bounds the error, and its sign tells you whether the estimate is high or low, since $f''<0$ (concave) means the tangent lies above the graph.

\begin{worked}
Estimate $\ln(1.05)$. With $f=\ln$, $a=1$: $f(1)=0$, $f'(1)=1$, so $L(1.05)=0.05$. Since $f''(x)=-1/x^{2}<0$, the estimate is an overshoot, and indeed $\ln(1.05)=0.0487902\ldots$. The bound $\tfrac12\abs{f''(\xi)}(0.05)^{2}\leq\tfrac12(1)(0.0025)=0.00125$ contains the actual error $0.00121$.

Estimate $\sin31^{\circ}$. Convert first: $31^{\circ}=\tfrac{\pi}{6}+\tfrac{\pi}{180}$, so $a=\tfrac{\pi}{6}$ and $h=\tfrac{\pi}{180}=0.0174533$. Then
\[\sin31^{\circ}\approx\sin\tfrac{\pi}{6}+\cos\tfrac{\pi}{6}\cdot\tfrac{\pi}{180}=0.5+0.8660254(0.0174533)=0.5151150,\]
against the true $0.5150381$. The estimate is high, as $\sin$ is concave on $(0,\pi)$.

Estimate $\sqrt[3]{28}$. The nearest exact point is $a=27$, where $f(27)=3$ and $f'(27)=\tfrac13(27)^{-2/3}=\tfrac{1}{27}$. So $\sqrt[3]{28}\approx3+\tfrac{1}{27}=3.037037$, against $3.036589$; again high, since $x^{1/3}$ is concave for $x>0$, and again the error $4.5\times10^{-4}$ is of the predicted size $\tfrac12\abs{f''(27)}=\tfrac12\cdot\tfrac{2}{9}\cdot27^{-5/3}=4.6\times10^{-4}$.
\end{worked}

\begin{trap}
Two failures, both common. First, choosing $a$ for arithmetic comfort rather than proximity: estimating $\sqrt{104}$ from $a=100$ is fine, but estimating $\sqrt{140}$ from $a=100$ gives $12.0$ against a true $11.832$, because $(x-a)^{2}=1600$ has made the quadratic term large. Second, linearising a trigonometric function without converting to radians. The formula $\dd(\sin x)=\cos x\dd x$ is false in degrees; in degrees the derivative of $\sin$ carries a factor $\pi/180$.
\end{trap}

\begin{seealsob}The differential versus the increment; error propagation; Taylor polynomials with explicit remainder.\end{seealsob}

Linear approximation stops being an approximation and starts being a lie at a rate controlled entirely by $f''$, and it is worth knowing the numbers. The error is $\tfrac12f''(\xi)(x-a)^{2}$, so halving the step quarters the error; a step ten times smaller gives an answer a hundred times better. Concretely, for $\sqrt{\cdot}$ at $a=100$: a step of $0.4$ gives an error of $2\times10^{-5}$, a step of $4$ gives $2\times10^{-3}$, a step of $40$ gives $0.19$, and a step of $400$ gives an answer, $12$, that is wrong in the first digit. The linear model degrades quadratically and then catastrophically, and the only defence is to keep $\abs{x-a}$ small, which means choosing $a$ well. When no nearby exact point exists, the correct response is not a larger step but a second-order term: $f(a+h)\approx f(a)+f'(a)h+\tfrac12f''(a)h^{2}$, which is the next line of the Taylor expansion and costs one more derivative.

\section{Error, and why relative error is the natural currency}

Approximation and error propagation are the same computation read in opposite directions. In approximation, $\dd x$ is a step you chose and $\dd y$ is the estimate you want. In error propagation, $\dd x$ is an uncertainty you were handed and $\dd y$ is the uncertainty you must report. Nothing changes except that absolute values appear, because a tolerance has no sign, and that the answer is usually wanted as a percentage rather than as a quantity.

\tech{Error propagation and relative error}{easy}
\cue A measured quantity carries a stated tolerance ($r=5\pm0.02$ cm, ``accurate to within $1\%$'') and you are asked for the induced tolerance in something computed from it.
\method Absolute error is $\abs{\dd y}=\abs{f'(x)}\abs{\dd x}$. Relative error is the quantity actually worth computing:
\[\left|\frac{\dd y}{y}\right|=\left|\frac{f'(x)}{f(x)}\right|\abs{\dd x}=\bigl|\bigl(\ln\abs{f}\bigr)'(x)\bigr|\,\abs{\dd x}.\]
For a power law $y=Cx^{n}$ this collapses to $\dfrac{\dd y}{y}=n\dfrac{\dd x}{x}$: relative error multiplies by the exponent and the constant $C$ is irrelevant. Take logs before differentiating whenever the formula is a product, quotient, or power, since $\ln$ turns all three into sums.

\begin{worked}
A sphere's radius is measured as $r=5\pm0.02$ cm. Find the induced uncertainty in the volume $V=\tfrac43\pi r^{3}$, absolutely and as a percentage.

By the power rule for relative error, $\dfrac{\dd V}{V}=3\dfrac{\dd r}{r}=3\cdot\dfrac{0.02}{5}=3(0.004)=0.012$, so $1.2\%$. Absolutely, $\dd V=4\pi r^{2}\dd r=4\pi(25)(0.02)=2\pi\approx6.283\ \mathrm{cm}^{3}$, against $V=\tfrac{500\pi}{3}\approx523.6\ \mathrm{cm}^{3}$; the ratio is $0.012$, as promised.

The same rule read backwards answers the design question. If the volume must be known to $0.5\%$, then $\dfrac{\dd r}{r}\leq\dfrac{0.005}{3}=0.00167$, so the radius must be measured to about $0.008$ cm. Tightening a volume tolerance costs three times as much precision in the radius.

The power rule handles negative and fractional exponents without change, and that is where it earns its keep. A pendulum has period $T=2\pi\sqrt{L/g}$, so $T=CL^{1/2}$ and $\dfrac{\dd T}{T}=\tfrac12\dfrac{\dd L}{L}$: a length known to $1\%$ gives a period known to $0.5\%$, and the exponent $\tfrac12$ is doing all the work. Inverted to measure gravity, $g=4\pi^{2}LT^{-2}$ gives
\[\frac{\dd g}{g}=\frac{\dd L}{L}-2\frac{\dd T}{T},\qquad \left|\frac{\dd g}{g}\right|\leq\left|\frac{\dd L}{L}\right|+2\left|\frac{\dd T}{T}\right|,\]
so timing error costs twice what length error costs. That asymmetry is invisible in the formula for $g$ and obvious in the formula for $\dd g/g$, which is the reason to compute the latter.
\end{worked}

\begin{trap}
Adding the \emph{absolute} errors of factors in a product. For $A=\ell w$ the correct combination is $\dfrac{\dd A}{A}=\dfrac{\dd\ell}{\ell}+\dfrac{\dd w}{w}$: relative errors add, absolute errors do not. Separately, dropping the absolute value, so that a decreasing function ($f'<0$) is reported as having a negative tolerance. A tolerance is a magnitude.
\end{trap}

\begin{seealsob}The total differential of a multivariable function; the differential versus the increment; logarithmic differentiation.\end{seealsob}

One measured input is the exceptional case. Real formulas take several, each with its own tolerance, and the linear-part idea extends verbatim: the increment of $f(x,y)$ is a linear function of the pair $(\dd x,\dd y)$, with the two partial derivatives as coefficients. The only genuinely new question is how to combine the two contributions, and the answer depends on a modelling assumption that the problem statement, not the mathematics, decides.

\tech{The total differential of a multivariable function}{medium}
\cue A quantity depends on several independently measured inputs, or a problem writes $\dd z$ for $z=f(x,y)$, or you need to know how a formula responds to simultaneous small changes in all of its arguments.
\method The linear part of the increment is
\[\dd z=\pdv{f}{x}\dd x+\pdv{f}{y}\dd y,\qquad \Delta z=\dd z+o\bigl(\abs{\dd x}+\abs{\dd y}\bigr).\]
For a guaranteed worst-case bound, use the triangle inequality: $\abs{\dd z}\leq\abs{f_x}\abs{\dd x}+\abs{f_y}\abs{\dd y}$, which assumes the errors conspire. For independent random errors with standard deviations $\sigma_x,\sigma_y$, combine in quadrature instead: $\sigma_z^{2}=f_x^{2}\sigma_x^{2}+f_y^{2}\sigma_y^{2}$.

\begin{worked}
Two resistors in parallel satisfy $\dfrac1R=\dfrac1{R_1}+\dfrac1{R_2}$. Differentiate the relation as it stands, without solving for $R$:
\[-\frac{\dd R}{R^{2}}=-\frac{\dd R_1}{R_1^{2}}-\frac{\dd R_2}{R_2^{2}},\qquad\text{so}\qquad \dd R=\frac{R^{2}}{R_1^{2}}\dd R_1+\frac{R^{2}}{R_2^{2}}\dd R_2.\]
Take $R_1=100\pm1$ and $R_2=400\pm4$ ohms. Then $\tfrac1R=\tfrac{1}{100}+\tfrac{1}{400}=\tfrac{1}{80}$, so $R=80$, and the worst-case bound is
\[\abs{\dd R}\leq\frac{6400}{10000}(1)+\frac{6400}{160000}(4)=0.64+0.16=0.80\ \Omega.\]
Direct computation confirms it: replacing $(100,400)$ by $(99,396)$ gives $R=79.2$ exactly, and by $(101,404)$ gives $R=80.8$. Note what the numbers say. Both inputs carry $1\%$ error and the output carries $0.80/80=1\%$; the parallel combination is homogeneous of degree one, so it inherits the common relative error unchanged. Under the quadrature rule the answer would instead be $\sqrt{0.64^{2}+0.16^{2}}=0.66\ \Omega$.

A second instance, where taking logarithms first shortens the work considerably. A cylinder has $r=2\pm0.01$ and $h=5\pm0.02$ cm, and $V=\pi r^{2}h$. Rather than computing $V_r$ and $V_h$, take logs: $\ln V=\ln\pi+2\ln r+\ln h$, so
\[\frac{\dd V}{V}=2\frac{\dd r}{r}+\frac{\dd h}{h}=2(0.005)+0.004=0.014,\]
a worst case of $1.4\%$. Since $V=20\pi=62.832$, that is $\dd V=0.880\ \mathrm{cm}^{3}$, and the exact worst case, $\pi(2.01)^{2}(5.02)-20\pi=0.884$, sits just above it as the second-order terms require. The logarithmic route generalises: for any monomial $z=Cx^{a}y^{b}$, the relative errors combine as $\abs{a}$ and $\abs{b}$ times the input relative errors, with no differentiation at all.
\end{worked}

\begin{trap}
Using quadrature when the problem asks for a guaranteed bound (``the volume cannot be wrong by more than''), or the linear sum when it asks for a standard deviation. These give different numbers, $0.80$ against $0.66$ above, and only the wording of the problem decides which is wanted. A second trap: computing $R=80$ first and then differentiating $R=\dfrac{R_1R_2}{R_1+R_2}$ by brute force. Differentiating the reciprocal relation in place is three symbols shorter and less error-prone, which is the whole reason to keep the relation implicit.
\end{trap}

\begin{seealsob}Error propagation and relative error; implicit differentiation in differential form; the exactness test, which asks when a given $M\dd x+N\dd y$ is itself a total differential.\end{seealsob}

\begin{keynote}[Worst case or standard deviation]
The two combination rules answer different questions and the problem statement, not the mathematics, chooses between them. ``The volume cannot be wrong by more than'' asks for a guaranteed bound and takes the linear sum of absolute contributions, which assumes the errors conspire to point the same way. ``What is the uncertainty in the volume'' in a laboratory context usually asks for a standard deviation and takes the quadrature sum, which assumes the errors are independent and random. The quadrature answer is always the smaller of the two, and for two equal contributions it is smaller by a factor $\sqrt2$. Reporting a quadrature figure where a bound was demanded understates the risk; reporting a bound where a standard deviation was demanded overstates the error by a factor that grows with the number of inputs.
\end{keynote}

\section{Differentials as relations between variables}

So far $\dd y$ has been computed from an explicit $y=f(x)$. The notation earns its keep when there is no explicit $f$. An implicit relation $F(x,y)=0$ constrains the pair $(\dd x,\dd y)$ by a single linear equation, and that equation is all you need: solving it for the ratio gives $y'$, and never solving the original relation for $y$ is the point.

\tech{Implicit differentiation in differential form}{easy}
\cue A relation $F(x,y)=0$ that cannot be solved for $y$ in closed form, or can be solved only at the cost of radicals and case splits, together with a request for $y'$, a tangent line, or the location of horizontal and vertical tangents.
\method Apply $\dd$ to both sides, treating $x$ and $y$ symmetrically and using the product and chain rules on every term. The result is linear in $\dd x$ and $\dd y$:
\[F_x\dd x+F_y\dd y=0\qquad\Longrightarrow\qquad \dv{y}{x}=-\frac{F_x}{F_y}.\]
Horizontal tangents are where $F_x=0$ and $F_y\neq0$; vertical tangents are where $F_y=0$ and $F_x\neq0$. Both are visible in the differential relation before any division, which is exactly why you should read it before dividing.

\begin{worked}
The folium of Descartes, $x^{3}+y^{3}=6xy$. Taking $\dd$ of both sides, with $\dd(xy)=y\dd x+x\dd y$:
\[3x^{2}\dd x+3y^{2}\dd y=6y\dd x+6x\dd y,\qquad (3x^{2}-6y)\dd x+(3y^{2}-6x)\dd y=0,\]
so $\dv{y}{x}=\dfrac{2y-x^{2}}{y^{2}-2x}$. At the point $(3,3)$, which is on the curve since $27+27=54=6\cdot9$, this gives $\dfrac{6-9}{9-6}=-1$: the tangent is $y=6-x$.

The vertical tangents come free. Set $y^{2}=2x$ and substitute into the curve: $x^{3}+y^{3}=6xy$ becomes $\tfrac{y^{6}}{8}+y^{3}=3y^{3}$, so $y^{3}=16$ and $y=2^{4/3}$, $x=2^{5/3}$. Check: $x^{3}+y^{3}=32+16=48$ and $6xy=6\cdot2^{3}=48$. The curve has a vertical tangent at $(2^{5/3},2^{4/3})\approx(3.1748,2.5198)$, a fact that would have been painful to extract from any explicit solution for $y$.

A cleaner instance, showing both tangent types at once: $x^{2}-xy+y^{2}=3$. Taking $\dd$ gives $(2x-y)\dd x+(2y-x)\dd y=0$, so $\dv{y}{x}=\dfrac{2x-y}{x-2y}$. At $(1,-1)$, which lies on the curve, the slope is $\dfrac{3}{3}=1$ and the tangent is $y=x-2$. Horizontal tangents need $2x-y=0$; substituting $y=2x$ into the curve gives $3x^{2}=3$, so they occur at $(1,2)$ and $(-1,-2)$. Vertical tangents need $x-2y=0$; substituting $x=2y$ gives $3y^{2}=3$, so they occur at $(2,1)$ and $(-2,-1)$. Four special points, all read off the differential relation, and the curve was never solved for $y$.
\end{worked}

\begin{trap}
Forgetting that $y$ depends on $x$ inside a product: $\dd(xy)=y\dd x+x\dd y$, never $y\dd x$ alone. And dividing by $F_y$ without recording where it vanishes; at those points the formula for $y'$ is not merely large, it does not exist, and the curve is doing something (a vertical tangent, a cusp, a self-intersection) that the quotient conceals. The folium passes through the origin twice, and there $F_x=F_y=0$: the differential relation degenerates to $0=0$ and carries no information at all.
\end{trap}

\begin{seealsob}The total differential of a multivariable function; the exactness test; differentials as substitution bookkeeping.\end{seealsob}

The implicit relation $F_x\dd x+F_y\dd y=0$ is a constraint linking two increments. A substitution is the same kind of object read as a change of variable rather than as a constraint, and this is where the notation is most often abused. The equation $\dd u=g'(x)\dd x$ is a \emph{statement about the substitution $u=g(x)$}, derived by applying the definition of the differential to $g$; it is not a fraction being cross-multiplied, and no cancellation is taking place.

\tech{Differentials as substitution bookkeeping}{core}
\cue Any integral or differential equation in which a substitution $u=g(x)$ is introduced and the $\dd x$ must be converted. In a first-order ODE this is exactly how $y=vx$, $u=x+y$, or $v=y^{1-n}$ transfers the equation to a new unknown.
\method Write $\dd u=g'(x)\dd x$, then \emph{solve for $\dd x$ before substituting anything}, so that no factor is silently absorbed:
\[\dd x=\frac{\dd u}{g'(x)},\qquad g'(x)\ \text{re-expressed in }u.\]
After the substitution, every occurrence of the old variable must be gone. If one survives, the substitution has not been carried out; it has been started.

\begin{worked}
Solve $\dv{y}{x}=(x+y)^{2}$ with $y(0)=0$. The right side depends on $x$ and $y$ only through $u=x+y$, so substitute. The differential relation for the substitution is $\dd u=\dd x+\dd y$, hence
\[\dv{u}{x}=1+\dv{y}{x}=1+u^{2}.\]
This is separable: $\dfrac{\dd u}{1+u^{2}}=\dd x$ gives $\arctan u=x+C$. The initial condition $y(0)=0$ means $u(0)=0$, so $C=0$ and $u=\tan x$, that is,
\[y=\tan x-x,\qquad y(\tfrac{\pi}{4})=1-\tfrac{\pi}{4}\approx0.214602.\]
A fourth-order Runge-Kutta integration from $x=0$ to $x=\pi/4$ returns $0.2146018$, confirming both the closed form and the branch. The solution exists only on $(-\tfrac\pi2,\tfrac\pi2)$, since $u=\tan x$ escapes to infinity at the endpoints.

Note precisely what was legitimate. We did not multiply anything by $\dd x$. We wrote the differential of the function $u=x+y$ of two variables using the total differential, $\dd u=\dd x+\dd y$, and then divided the resulting identity between real numbers by the nonzero number $\dd x$.
\end{worked}

\begin{trap}
Two errors, and both are silent. In a definite integral, converting the integrand but not the limits: $\int_0^{1}x\eu^{x^{2}}\dd x$ with $u=x^{2}$ becomes $\tfrac12\int_0^{1}\eu^{u}\dd u$ because $u$ also runs from $0$ to $1$ here, but $\int_1^{2}$ would not be so forgiving. In an ODE, converting $\dd y$ while leaving a stray $y$ behind, which produces an equation in two unknowns that then gets ``solved'' by treating one of them as a constant.
\end{trap}

\begin{seealsob}Substitution $u=ax+by+c$; the substitution $y=vx$ for homogeneous equations; recognising a differential as $\dd(\text{something})$.\end{seealsob}

The same bookkeeping runs an integral, and it is worth seeing the two side by side because the discipline is identical. In $\displaystyle\int_1^{2}\frac{x\dd x}{1+x^{2}}$, put $u=1+x^{2}$. Then $\dd u=2x\dd x$, so $x\dd x=\tfrac12\dd u$, and the limits move with the variable: $x=1$ gives $u=2$ and $x=2$ gives $u=5$. Hence the integral is $\tfrac12\int_2^{5}\dfrac{\dd u}{u}=\tfrac12\ln\tfrac52=0.4581454$. Nothing was cancelled and nothing was multiplied by $\dd x$; the differential relation $\dd u=2x\dd x$ was solved for the group $x\dd x$ that actually appears, which is the general procedure. When the group you need is not present in the integrand, the substitution is the wrong one, and no amount of algebra on $\dd x$ will fix it.

The last technique in the chapter reverses the direction of every previous one. Instead of computing the differential of a known expression, you are handed a combination of $\dd x$ and $\dd y$ and must recognise it as the differential of something. This is the single highest-yield pattern-matching skill in first-order differential equations, because a recognised combination turns an equation into one integration.

\tech[Recognising a differential as d of something]{Recognising a differential as $\dd(\text{something})$}{medium}
\cue An equation containing the combination $x\dd y+y\dd x$, or $x\dd y-y\dd x$, or $tP'+P$, or any side that is visibly the derivative of a named expression. A hint that this is intended: the equation is not separable, not linear, and yet looks tidy.
\method Learn the integrable combinations by heart. Each is a product or quotient rule read backwards:
\[\dd(xy)=x\dd y+y\dd x,\qquad \dd\!\left(\frac yx\right)=\frac{x\dd y-y\dd x}{x^{2}},\]
\[\dd\!\left(\arctan\frac yx\right)=\frac{x\dd y-y\dd x}{x^{2}+y^{2}},\qquad \dd\!\left(\ln\frac yx\right)=\frac{x\dd y-y\dd x}{xy}.\]
The last three all begin with $x\dd y-y\dd x$; what distinguishes them is what you divide by, so identify that combination first and then look at the denominator you have been given. Two more are worth adding to the list because they turn up in exact equations: $\dd\bigl(\tfrac12\ln(x^{2}+y^{2})\bigr)=\dfrac{x\dd x+y\dd y}{x^{2}+y^{2}}$, and $\dd\bigl(x^{a}y^{b}\bigr)=x^{a-1}y^{b-1}\bigl(ay\dd x+bx\dd y\bigr)$, which is the general monomial and which contains $\dd(xy)$ as the case $a=b=1$.

\begin{worked}
Solve $2y\dv{y}{x}=\sin x+x\cos x$ with $y(0)=0$. Neither side needs integrating in the usual sense: the left is $\dvop{x}(y^{2})$ and the right is $\dvop{x}(x\sin x)$, since $\dd(x\sin x)=\sin x\dd x+x\cos x\dd x$. Two functions with equal derivatives differ by a constant, and $y(0)=0$ makes it zero:
\[y^{2}=x\sin x.\]
A second instance, where the packaging is a quotient. Solve $x\dv{y}{x}-y=x^{2}+y^{2}$. Divide by $x^{2}$ to expose the combination:
\[\frac{x\dd y-y\dd x}{x^{2}}=\left(1+\frac{y^{2}}{x^{2}}\right)\dd x\quad\Longrightarrow\quad \frac{\dd(y/x)}{1+(y/x)^{2}}=\dd x,\]
so $\arctan\dfrac yx=x+C$ and $y=x\tan(x+C)$. With $y(\tfrac\pi2)=\tfrac\pi2$ we need $\tan(\tfrac\pi2+C)=1$, giving $C=-\tfrac\pi4$ and $y=x\tan\!\left(x-\tfrac\pi4\right)$; then $y(\tfrac\pi4)=0$ and $y(1.2)=0.5281360$, which matches Runge-Kutta to seven figures.

A third, the one that catches people: $t\dv{A}{t}-A=t^{3}$. Dividing by $t^{2}$ gives $\dvop{t}\!\left(\dfrac At\right)=t$, so $\dfrac At=\dfrac{t^{2}}{2}+C$ and $A=\dfrac{t^{3}}{2}+Ct$.
\end{worked}

\begin{trap}
Looking only for the product rule. The three most useful combinations are quotient-rule packagings, and the giveaway is $x\dd y-y\dd x$ with a minus sign, which no product rule produces. Dividing by the wrong power is the other half of the same mistake: $x\dd y-y\dd x$ over $x^{2}$ gives $\dd(y/x)$, over $y^{2}$ gives $-\dd(x/y)$, and over $xy$ gives $\dd\ln(y/x)$. Choose the one whose denominator the equation actually hands you rather than the one you remember best.
\end{trap}

\begin{seealsob}Differentials as substitution bookkeeping; the exactness test; the integrating-factor machine, which manufactures the combination when it is not already present.\end{seealsob}

The combinations are few enough to tabulate, and the table is worth memorising as a unit rather than one entry at a time, because in practice you recognise the numerator first and then hunt through the possible denominators.

\begin{center}
\footnotesize
\setlength{\tabcolsep}{5pt}
\renewcommand{\arraystretch}{1.45}
\begin{tabular}{@{}lll@{}}
\toprule
\mh{Combination} & \mh{Divide by} & \mh{Equals} \\
\midrule
$x\dd y+y\dd x$ & $1$ & $\dd(xy)$ \\
$x\dd y+y\dd x$ & $xy$ & $\dd\ln\abs{xy}$ \\
$x\dd y-y\dd x$ & $x^{2}$ & $\dd\bigl(\tfrac yx\bigr)$ \\
$x\dd y-y\dd x$ & $y^{2}$ & $-\dd\bigl(\tfrac xy\bigr)$ \\
$x\dd y-y\dd x$ & $xy$ & $\dd\ln\abs{\tfrac yx}$ \\
$x\dd y-y\dd x$ & $x^{2}+y^{2}$ & $\dd\arctan\tfrac yx$ \\
$x\dd x+y\dd y$ & $x^{2}+y^{2}$ & $\dd\bigl(\tfrac12\ln(x^{2}+y^{2})\bigr)$ \\
$x\dd x+y\dd y$ & $\sqrt{x^{2}+y^{2}}$ & $\dd\sqrt{x^{2}+y^{2}}$ \\
$ay\dd x+bx\dd y$ & $x^{1-a}y^{1-b}$ & $\dd\bigl(x^{a}y^{b}\bigr)$ \\
\bottomrule
\end{tabular}
\end{center}

The last row subsumes the first and is the one to reach for when the coefficients of $y\dd x$ and $x\dd y$ are different constants: an equation containing $3y\dd x+2x\dd y$ is asking to be multiplied by $x^{2}y$, after which it is $\dd(x^{3}y^{2})$. Reading the table in that direction, from a mismatched pair of coefficients to the monomial factor that repairs them, is exactly the search that the integrating-factor method automates for exact equations. Recognising the combination by eye is faster whenever it works, and it works often enough to be the first thing tried.

\section{Which manipulations are legitimate}

It is worth collecting the verdicts, because the rules are simple and the folklore around them is not. Writing $\dd y=f'(x)\dd x$ is a definition and always legal. Dividing that identity by the nonzero real number $\dd x$ is legal, and gives $\dd y/\dd x=f'(x)$. Chaining is legal: if $y=f(u)$ and $u=g(x)$, then $\dd y=f'(u)\dd u=f'(g(x))g'(x)\dd x$, and the chain rule is literally the statement that these two computations of $\dd y$ agree. Writing $\int f(g(x))g'(x)\dd x=\int f(u)\dd u$ is legal and is the substitution theorem, not a cancellation.

What is not legal is any move that requires $\dd x$ to be nonzero and infinitely small at the same time, and any move that treats $\dd y$ and $\dd x$ as separately meaningful outside a ratio, an integral, or a $1$-form. Separating variables, written as $h(y)\dd y=g(x)\dd x$, is a compressed instruction to integrate $h(y(x))y'(x)$ against $x$ and apply the substitution theorem; it is legal, and the legality is a theorem about integrals rather than an algebraic identity about symbols. Finally, $\dd(\dd y)$ is not an ordinary second differential: $\dd^{2}y/\dd x^{2}$ is a single symbol for $(y')'$, and treating the numerator and denominator as separate powers is the one place where Leibniz notation genuinely misleads, since $\ddv{y}{x}\neq\left(\dv{y}{x}\right)^{2}$ and the correct chain rule for second derivatives carries an extra term.

There is one more habit worth naming, because it is what converts the notation from a hazard into a tool. When a problem hands you a relation between quantities, differentiate the relation \emph{in place} rather than solving it first. The parallel-resistor computation above is the model: differentiating $\tfrac1R=\tfrac1{R_1}+\tfrac1{R_2}$ takes one line, while solving for $R$ and then differentiating the quotient takes five and introduces three opportunities for error. The same applies to implicit curves, to thermodynamic relations such as $PV=nRT$ (whose differential $\dfrac{\dd P}{P}+\dfrac{\dd V}{V}=\dfrac{\dd T}{T}$ is the whole of the gas law's error analysis), and to any formula defined by a constraint rather than by an explicit expression. The differential of a relation is a relation between differentials, and it is almost always simpler than the thing it came from.

\begin{keynote}[The one-line test]
Before any manipulation of $\dd x$, ask: which of the three readings am I using, and is the move licensed there? Approximation and error work is reading (i), finite increments, where every equation is an identity between real numbers. Substitution and separation are reading (ii), where the only rule is $u=g(x)\Rightarrow\dd u=g'(x)\dd x$ and the justification is the substitution theorem. Exactness and path independence are reading (iii). If a manipulation is licensed in none of the three, it is not licensed.
\end{keynote}

\section{Recognition matrix}
\begin{recog}{@{}p{0.30\textwidth}p{0.36\textwidth}p{0.28\textwidth}@{}}
\rowcolor{mist}\mh{If you see} & \mh{Reach for} & \mh{If that fails} \\
\midrule
\endhead
``Approximately how much does $y$ change'' & $\dd y=f'(x)\dd x$ with $\dd x=\Delta x$ & Second-order term $\tfrac12f''(\dd x)^{2}$ \\
``Estimate $f(a+h)$'', $a$ clean & $L(x)=f(a)+f'(a)(x-a)$ & Taylor to degree two \\
Concavity known, error sign wanted & Sign of $f''$: concave means overshoot & Lagrange remainder at $\xi$ \\
A tolerance $\pm\delta$ on one input & $\abs{\dd y}=\abs{f'}\delta$; then divide by $y$ & Endpoint evaluation of $f$ \\
A product, quotient, or power & Log first: $\dd y/y=(\ln\abs f)'\dd x$ & Term-by-term relative errors \\
$y=Cx^{n}$ with $x$ uncertain & $\dfrac{\dd y}{y}=n\dfrac{\dd x}{x}$; $C$ irrelevant & Direct differentiation \\
Several measured inputs, worst case & $\abs{\dd z}\leq\abs{f_x}\abs{\dd x}+\abs{f_y}\abs{\dd y}$ & Evaluate at all corner points \\
Several inputs, random and independent & Quadrature: $\sigma_z^{2}=f_x^{2}\sigma_x^{2}+f_y^{2}\sigma_y^{2}$ & Read the wording again \\
$F(x,y)=0$, unsolvable for $y$ & $F_x\dd x+F_y\dd y=0$, then divide & Parametrise the curve \\
Vertical tangent wanted & Solve $F_y=0$ jointly with $F=0$ & Swap roles: find $\dd x/\dd y=0$ \\
$\dd(xy)$ appearing as $y\dd x$ & Product rule: restore $x\dd y$ & Rewrite $F$ before differentiating \\
A substitution $u=g(x)$ to be carried out & $\dd u=g'\dd x$; solve for $\dd x$ first & Re-express $g'$ in terms of $u$ \\
$x\dd y+y\dd x$ & $\dd(xy)$; integrate once & Look for an integrating factor \\
$x\dd y-y\dd x$ over $x^{2}$ & $\dd(y/x)$; substitute $v=y/x$ & Over $xy$: $\dd\ln(y/x)$ \\
$x\dd y-y\dd x$ over $x^{2}+y^{2}$ & $\dd\arctan(y/x)$ & Polar coordinates \\
$tA'-A$ or $tP'+P$ & Divide by $t^{2}$, or read $(tP)'$ & Standard integrating factor \\
$\dd^{2}y/\dd x^{2}$ treated as a square & It is not; use the chain rule with its extra term & Write $p=y'$ and reduce order \\
\end{recog}
